% Chapter 4: panel buttons and LEDs.
\chapter{Panel buttons and LEDs}\label{ch:panel}
\index{buttons}\index{LEDs}

\section{How the buttons work}\label{sec:latching}
\index{buttons!latching}
The four panel buttons are \textbf{latching} switches. Push one and it stays
pushed in; push it again and it pops back out. The machine does not care
which way a button is resting. \textbf{Every push, in or out, counts as one
click}, and one click does one thing.

\begin{itemize}
  \item After a wind a button may still be sitting pushed in. That is normal.
        You do not need to pop it out before the next wind.
  \item A button that is already pushed in when you switch the power on does
        nothing. The machine ignores button positions at power-up. Wait about
        one second after power-up before you push a button.
  \item Very fast double pushes may count as one: the machine ignores
        contact chatter shorter than about 1/60 of a second.
  \item While the PC program has the machine in Test mode, the panel cannot
        start anything (Section~\ref{sec:testmode}).
\end{itemize}

\section{Button map}\label{sec:buttonmap}
\index{button map}

\begin{center}
\renewcommand{\arraystretch}{1.5}
\begin{tabularx}{\textwidth}{@{}l>{\raggedright\arraybackslash}X>{\raggedright\arraybackslash}X@{}}
\toprule
\textbf{Button} & \textbf{At the idle screen} & \textbf{During a wind} \\
\midrule
\btn{1} & Start a \textbf{30 yard} wind. & \textbf{Stop} the wind. \\
\btn{2} & Start a \textbf{90 yard} wind. & \textbf{Stop} the wind. \\
\btn{3} & Nothing. & \textbf{Pause}; push again to \textbf{resume}. \\
\btn{4} & \textbf{Manual}: motor on; push again for motor off. &
  While \emph{paused}: motor on/off to nudge the paper by hand
  (stops by itself at the target length). While winding: nothing. \\
\midrule
Any button & Clears an \shows{UNDER TARGET} alarm \emph{and} does its normal job. & \\
\bottomrule
\end{tabularx}
\end{center}

\begin{note}[Tip]
To stop quickly, push \emph{either} \btn{1} or \btn{2}; it does not matter
which one started the wind.
\end{note}

\section{LEDs}\label{sec:leds}
The four LEDs above the buttons tell you the state at a glance.

\begin{center}
\renewcommand{\arraystretch}{1.4}
\begin{tabularx}{\textwidth}{@{}l>{\raggedright\arraybackslash}X@{}}
\toprule
\textbf{LEDs lit} & \textbf{Meaning} \\
\midrule
\led{1} and \led{3} & Idle, ready to wind. \\
\led{4} only & Manual mode: the motor is on. \\
None, then 1, 1--2, 1--3, all four & Winding: a bar graph that fills in quarters
  as the wind goes on. \\
1--3, then all four & Trimming: all four light once the length has passed the target. \\
\led{2} and \led{4} & Test mode (PC program has control). \\
Fault LED (where fitted) & The controller found an internal fault. Switch off,
  then on; if it comes back, call [SUPPORT CONTACT]. \\
\bottomrule
\end{tabularx}
\end{center}

\section{Buzzer}
The buzzer beeps when a wind finishes, and sounds an alarm when a roll comes
up short. See Chapter~\ref{ch:beeps} for every pattern.
